\Needspace{18\baselineskip}
\begin{proposition}\label{prop:long-perron-cutoff-type-stability}
For $u\in\PSH(X,\theta)$ with $u\leq\phi$, put
\[
 u_T=\max\{u,\phi-T\},\qquad
 L_s^{u,T}=G_{s/T}(\phi,u_T),\qquad T>0,\quad0\leq s\leq T.
\]
For fixed $s$, the functions $L_s^{u,T}$ increase with $T\geq s$ and satisfy $\phi-s\leq L_s^{u,T}\leq\phi$. Their canonical limit
\[
 R_s^u=\left(\sup_{T>s}L_s^{u,T}\right)^*
\]
is $\theta$-psh. If $u,v\leq\phi$ and $u\leq v+C$ for $C\geq0$, then
\[
 L_s^{u,T}\leq L_s^{v,T}+\frac{Cs}{T},\qquad R_s^u\leq R_s^v.
\]
In particular $u\sim v$ implies $R_s^u=R_s^v$ for every $s\geq0$. For arbitrary $u\preceq\phi$, subtracting any constant that puts $u$ below $\phi$ defines the same profiles. They increase under the singularity preorder.
\end{proposition}
\begin{proof}
The endpoints $u_T$ have bounded type relative to $\phi$, so their Perron segments exist by \cref{thm:relative-full-mass-perron-geodesic}. The family $q_s=\max\{u,\phi-s\}$ on $[0,T]$ is a subgeodesic: its complexification is the maximum of the constant potential $u$ and the affine potential $\phi-\Real\zeta$. It has the required endpoints, since $q_0=\phi$, $q_T=u_T$ and $|q_s-q_t|\leq|s-t|$ on their common finite locus gives $L^1$ endpoint convergence. Hence $L_s^{u,T}\geq q_s\geq\phi-s$. Convexity and $u_T\leq\phi$ give the upper bound.

If $T'>T$, the same competitor gives $L_T^{u,T'}\geq u_T$. Restriction of the segment to $[0,T]$, followed by endpoint order, gives $L_s^{u,T'}\geq L_s^{u,T}$ for $s\leq T$. The upper envelope theorem and the two bounds therefore give a noncollapsed $\theta$-psh canonical limit for every fixed $s$. Its supremum may be taken over any cofinal tail of $T>s$; at $s=0$ all functions and the limit are $\phi$.

The terminal inequality implies
\[
 u_T\leq\max\{v+C,\phi-T\}\leq v_T+C.
\]
Endpoint order and the affine endpoint-shift rule of \cref{prop:perron-geodesic-order-and-limits} give the stated $Cs/T$ bound. For fixed $s$ and $A>s$, take the canonical supremum over $T\geq A$ to obtain
\[
 R_s^u\leq R_s^v+\frac{Cs}{A}.
\]
Let $A\to\infty$. Applying this argument in both directions proves bounded-type invariance. Any two permitted normalizations differ by a constant, so they give the same profiles. Finally $u\preceq v$ means $u\leq v+C$ for some finite $C$; enlarge it to be nonnegative and apply the same comparison to normalized potentials. Bounded-shift invariance removes the choice of normalization.
\end{proof}
\index{Perron envelope!long-segment cutoff stability}
